%basic packages
\usepackage[utf8]{inputenc}
\usepackage[T1]{fontenc}
\usepackage{graphicx}
\usepackage[margin=1in]{geometry}
\usepackage[usenames,dvipsnames]{xcolor}
\usepackage{microtype}
\usepackage[skip=5pt, indent=40pt]{parskip}


%math
\usepackage{amsmath, amsthm, amsfonts, amssymb, mathtools}
\usepackage{mathrsfs}
\usepackage{cancel}
\usepackage{siunitx} %phyjsicsssss
\usepackage{bbm} %mathbb for numbers
\usepackage{tikz-cd} % commutative diagrams
\usepackage{eucal}
\makeatletter
\renewcommand*\env@matrix[1][c]{\hskip -\arraycolsep
  \let\@ifnextchar\new@ifnextchar
  \array{*\c@MaxMatrixCols #1}}
\makeatother %matrix realignment

%misc
\usepackage{float}
\usepackage[hyphens]{url}
%\definecolor{page}{HTML}{242526}
%\pagecolor{page}
\usepackage{booktabs} %the \toprule and \bottomrule thick lines on tables

\usepackage{hyperref}
\definecolor{aqua}{HTML}{00C5FF}
%\hypersetup{
%    colorlinks,
%    linkcolor={aqua},
%    urlcolor={aqua},
%    citecolor={red}
%}
\usepackage{cleveref}

%my commands
\DeclarePairedDelimiter\bra{\langle}{\rvert} %Bra
\DeclarePairedDelimiter\ket{\lvert}{\rangle} %Ket
\DeclarePairedDelimiterX\braket[2]{\langle}{\rangle}{#1\,\delimsize\vert\,\mathopen{}#2} %Bra-ket
\DeclareMathOperator{\diag}{diag}
\DeclareMathOperator*{\argmax}{arg\,max}
\DeclareMathOperator*{\argmin}{arg\,min}
\DeclareMathOperator{\Hom}{Hom}
\DeclareMathOperator{\Obj}{Obj}
\DeclareMathOperator{\End}{End}
\DeclareMathOperator{\Aut}{Aut}
\DeclareMathOperator{\coker}{coker}
\DeclareMathOperator{\spec}{Spec}
\DeclareMathOperator{\im}{im}
\newcommand{\Rmod}{\mathcal{M}\mathrm{od}}
\newcommand{\Vect}{\mathcal{V}\mathrm{ect}}
\newcommand{\Alg}{\mathcal{A}\mathrm{lg}}
\newcommand{\tildesim}{\mathbin{\sim}}
\renewcommand{\mod}[1]{\;(\text{mod}\;#1)}
\DeclareMathOperator{\Tr}{Tr}

\newcommand{\Set}{\mathcal{S}\mathrm{et}}
\newcommand{\Grp}{\mathcal{G}\mathrm{rp}}
\newcommand{\Top}{\mathcal{T}\mathrm{op}}
\newcommand{\Ring}{\mathcal{R}\mathrm{ing}}
\newcommand{\Cat}{\mathcal{C}\mathrm{at}}
\newcommand{\Ab}{\mathcal{A}\mathrm{b}}
\DeclareMathOperator{\Nat}{Nat}
\DeclareMathOperator{\Fun}{Fun}
\DeclareMathOperator{\Map}{Map}
\DeclareMathOperator{\Lan}{Lan}
\DeclareMathOperator{\Ran}{Ran}
\newcommand{\sSet}{\mathrm{s}\mathcal{S}\mathrm{et}}
\renewcommand{\Im}{\operatorname{Im}}
\DeclareMathOperator{\Ker}{Ker}
\DeclareMathOperator{\Coker}{Coker}
\DeclareMathOperator{\Coim}{Coim}
\DeclareMathOperator*{\hocolim}{hocolim}
\DeclareMathOperator*{\holim}{holim}
%theorems
\usepackage{thmtools}
\usepackage{tikz}
\usepackage{tikz-cd}
\usepackage[framemethod=TikZ]{mdframed}
\mdfsetup{skipabove=1em,skipbelow=0em, innertopmargin=5pt, innerbottommargin=6pt}

%\declaretheoremstyle[headfont=\bfseries, bodyfont=\normalfont, qed={$\heartsuit$}]{definition}
%\declaretheorem[numberwithin=chapter, style=definition, name=Definition]{definition}
%\newtheorem{theorem}{Theorem}
%\newtheorem{corollary}{Corollary}
%\newtheorem{lemma}{Lemma}
%\newtheorem{proposition}{Proposition}

%theorem styles
%\declaretheoremstyle[headfont=\bfseries, bodyfont=\normalfont, mdframed={bottomline=false, topline=false, rightline=false, innerrightmargin=0pt}, qed=\(\blacksquare\)]{proofline}
%\declaretheoremstyle[headfont=\bfseries, bodyfont=\normalfont, mdframed={bottomline=false, topline=false, rightline=false, innerrightmargin=0pt}, qed=\qedsymbol]{egline}

%solution environment
%\declaretheorem[numbered=no, style=egline, name=Solution]{setsolution}
%\newenvironment{solution}[1][]{\vspace{-10pt}\begin{setsolution}}{\end{setsolution}}
%proof environment
%\declaretheorem[numbered=no, style=proofline, name=Proof]{replacementproof}
%\renewenvironment{proof}[1][\proofname]{\begin{replacementproof}}{\end{replacementproof}}

% Side Indented Theorems - https://tex.stackexchange.com/questions/429339/shifting-newtheorem
%\newtheoremstyle{side}{}{}{\advance\leftskip3cm\relax\itshape\normalfont}{-4pt}
%{\bfseries}{}{0pt}{
%\makebox[0pt][r]{
%  \smash{\parbox[t]{2.5cm}{\raggedright\thmname{#1}.
%  \thmnote{\newline(#3)}}}
%  \hspace{10.1pt}}}

%\theoremstyle{side}
%\newtheorem*{note}{Note}
%\newtheorem*{intuition}{Intuition}
%\newtheorem{claim}{Claim}
%\newtheorem*{prev}{As Previously Seen}

%\theoremstyle{definition}
%\newtheorem*{remark}{Remark}
%\newtheorem*{example}{Example}
%\newtheorem*{notation}{Notation}
%\newtheorem{exercise}{Exercise}


%fancy headers
\usepackage{fancyhdr}
\pagestyle{fancy}
\fancyhead{}\fancyfoot{}
\fancyfoot[R]{\thepage}
\fancyfoot[C]{\leftmark}
\renewcommand{\headrulewidth}{0pt}


\usepackage{pgfplots}

\author{Grant Talbert}



%%% darkmode because my pdf viewer on windows is traaaaaash


%\definecolor{darkbg}{rgb}{0.12, 0.12, 0.12}
%\definecolor{lightfg}{rgb}{0.84, 0.84, 0.84}
%\definecolor{accentblue}{rgb}{0.20, 0.63, 1.0}
\definecolor{darkbg}{RGB}{54, 57, 63}
\definecolor{lightfg}{RGB}{220, 221, 222}
%\definecolor{accentblue}{RGB}{88, 101, 242}
\definecolor{accentblue}{RGB}{133, 143, 255}
\pagecolor{darkbg}
\color{lightfg}
\usepackage{sectsty}
\allsectionsfont{\color{accentblue}}
\usepackage{hyperref}
\hypersetup{
    colorlinks=true,
    linkcolor=accentblue,
    filecolor=accentblue,      
    urlcolor=accentblue,
    citecolor=accentblue,
}

%%%%% theoremstyles
\declaretheoremstyle[headfont=\bfseries\color{accentblue}, bodyfont=\normalfont, qed={\color{accentblue}$\heartsuit$}]{definition}
\declaretheoremstyle[headfont=\bfseries\color{accentblue}, bodyfont=\normalfont\color{lightfg}, mdframed={
	bottomline=false, topline=false, rightline=false, innerrightmargin=0pt, linecolor=accentblue, backgroundcolor=darkbg, fontcolor={lightfg}, innertopmargin=0pt, innerbottommargin=0pt
}, qed={\color{accentblue}$\blacksquare$}]{proofline}
\declaretheoremstyle[headfont=\bfseries\color{accentblue}, bodyfont=\normalfont\color{lightfg}, mdframed={
	bottomline=false, topline=false, rightline=false, innerrightmargin=0pt, linecolor=accentblue, backgroundcolor=darkbg, innertopmargin=0pt, innerbottommargin=0pt, fontcolor={lightfg}
}, qed={\color{accentblue}\qedsymbol}]{egline}
\declaretheoremstyle[headfont=\bfseries\color{accentblue}, bodyfont=\normalfont]{dfn}
\declaretheoremstyle[headfont=\bfseries\color{accentblue}, bodyfont=\itshape]{thm}


\declaretheorem[numberwithin=chapter, style=definition, name=Definition]{definition}
\declaretheorem[numbered=no, style=dfn, name=Remark]{remark}
\declaretheorem[numbered=no, style=dfn, name=Example]{example}
\declaretheorem[numbered=no, style=dfn, name=Notation]{notation}
\declaretheorem[numbered=no, style=dfn, name=Exercise]{exercise}
\declaretheorem[numberwithin=chapter, style=thm, name=Theorem]{theorem}
\declaretheorem[numberwithin=chapter, style=thm, name=Corollary]{corollary}
\declaretheorem[numberwithin=chapter, style=thm, name=Lemma]{lemma}
\declaretheorem[numberwithin=chapter, style=thm, name=Proposition]{proposition}

%solution environment
\declaretheorem[numbered=no, style=egline, name=Solution]{setsolution}
\newenvironment{solution}[1][]{\vspace{-10pt}\begin{setsolution}}{\end{setsolution}}
%proof environment
\declaretheorem[numbered=no, style=proofline, name=Proof]{replacementproof}
\renewenvironment{proof}[1][\proofname]{\begin{replacementproof}}{\end{replacementproof}}

% Side Indented Theorems - https://tex.stackexchange.com/questions/429339/shifting-newtheorem
\newtheoremstyle{side}{}{}{\advance\leftskip3cm\relax\itshape\normalfont}{-4pt}
{\bfseries\color{accentblue}}{}{0pt}{
\makebox[0pt][r]{
  \smash{\parbox[t]{2.5cm}{\raggedright\thmname{#1}.
  \thmnote{\newline(#3)}}}
  \hspace{10.1pt}}}

\theoremstyle{side}
\newtheorem*{note}{Note}
\newtheorem*{intuition}{Intuition}
\newtheorem{claim}{Claim}
\newtheorem*{prev}{As Previously Seen}
